% Aportación separada, 9 de septiembre de 2026.
% RESULTADO_RECUPERADO: construcción sectorial del propietario 08_ckm_cp.tex.
% FORMALIZACION_REUNIDA: descomposición de paridad y compatibilidad representacional.
% Insertar antes de la carta racional CKM; conservar el desarrollo posterior.
% La procedencia y el alcance exacto están en 08b_procedencia_angular.md.

\subsection{Incidence sectorielle et construction du lecteur angulaire}
\label{sec:ii09-construccion-sectorial}

La construction angulaire utilise les représentations d’action et de clôture
obtenues au préalable à partir d’APP--TRIT--TPK. L’angle direct, le
contre-angle et le défaut torsionnel conservent leurs domaines propres.
Le lecteur des quarks combine ces trois coordonnées au moyen de l’incidence
sectorielle ; son évaluation précède les amplitudes complexes et la
comparaison métrologique. L’inclusion hexadique--octadique intervient dans
le troisième secteur à travers la différence de leurs cardinaux.

% ARQUITECTURA_AUTORAL_PREEXISTENTE: semicuadrado bidireccional de N42.
% FORMALIZACION_NUEVA: grupoide de pares y ponderacion por estabilizadores.
\subsubsection{Incidence marquée et normalisation bidirectionnelle}
\label{sec:ii09-normalizacion-bidireccional}

Les invariants du lecteur conservent la nature de l'objet dont ils
proviennent. L'intersection marquée $B_K=H_e\cap P_\pi$ a pour cardinal
$b=4$ ; l'entier $q=7$ est la coordonnée du pivot $p_\varphi$ dans la
numérotation orientée. L'ordre de détermination du système de Steiner
est $p=5$. Les supports hexadique et octadique ont pour cardinaux $h_6=6$
et $o_8=8$, et l'alphabet d'orientation tritique a pour cardinal $t=3$.
Un cardinal de support et une coordonnée de pivot sont utilisés par
leurs applications respectives, même lorsqu'ils sont tous deux publiés comme entiers.

La normalisation de l'autocroisement pentagonal admet la
réalisation explicite suivante. Soit $P$ le support de cinq positions de la clôture
et soit $X=P\times P$ son espace de paires ordonnées. L'inversion
bidirectionnelle sur cet espace est
\[
 \jmath:X\longrightarrow X,\qquad \jmath(x,y)=(y,x),
 \qquad\jmath^2=I.
\]
L'information de stabilisateur de chaque orbite est conservée. Pour une
orbite $[z]$, on définit son poids par
$1/|\operatorname{Stab}_{C_2}(z)|$. La somme de ces poids est la
cardinalité du groupoïde d'action $X/\!/C_2$.

\begin{proposition}[Mesure orbitale de l'autocroisement bidirectionnel]
\label{prop:ii09-semicuadrado-orbital}
Pour un support fini de cardinal $p$, l'inversion des paires ordonnées
détermine
\[
 |(P\times P)/\!/C_2|
 =\sum_{[z]\in(P\times P)/C_2}
      \frac1{|\operatorname{Stab}_{C_2}(z)|}
 =\frac{p^2}{2}.
\]
Pour $p=5$, l'autocroisement possède dix orbites libres et cinq orbites
fixes, et sa mesure orbitale est $10+5/2=25/2$.
\end{proposition}
\begin{proof}
Chaque paire diagonale $(x,x)$ forme une orbite fixe, dont le stabilisateur
est d'ordre deux ; il y en a $p$. Les $p(p-1)$ paires restantes se
regroupent deux par deux, avec un stabilisateur trivial. Leur contribution est
donc $p(p-1)/2$, et celle de la diagonale est $p/2$. La somme vaut
$p^2/2$. De manière équivalente, l'identité orbite--stabilisateur donne
$1/|\operatorname{Stab}(z)|=|[z]|/2$ ; la sommation sur toutes les orbites
produit $|P\times P|/2$.
\end{proof}

La mesure conserve explicitement la contribution des points fixes.
L'ensemble ordinaire des orbites a pour cardinal $p(p+1)/2$, qui vaut
quinze dans la clôture pentagonale. La mesure orbitale et ce cardinal sont
deux lectures déterminées de la même inversion. La normalisation
bidirectionnelle du lecteur utilise la première : les dix orbites libres
pèsent une unité et les cinq diagonales pèsent une demi-unité.

Cette réalisation précise le facteur $p^2/2$ employé dans le système
sectoriel suivant. L'affectation de cette contribution au secteur $23$,
son orientation négative et les deux autres règles d'incidence restent
spécifiées par le système sectoriel complet. La démonstration établit la
normalisation de l'autocroisement et maintient distincte la sélection des
fonctionnelles angulaires.


% FORMALIZACION_NUEVA de la elevacion de N39, compatible con (rho9,q9),
% T_d y Upd_t; arquitectura de residuo/cociente y orientacion preexistente.
% Fragmento de investigacion separado. No selecciona las reglas CKM N42.
\subsubsection{Transport de retenue arithmétique et holonomie entière du croisement arithmétique}
\label{subsec:transporte-entero-cruce}

La lecture modulaire du croisement admet une réalisation entière qui maintient
l’information accumulée lors du franchissement du bord. On développe ici cette
réalisation au moyen des quotients de transport. Son domaine est une coordinatisation
relevée du support APP ; l’assignation ultérieure d’une région à un secteur
angulaire conserve sa propre application d’incidence.

\paragraph{Coordinatisation nonadique et transport relevé.}
On emploie la décomposition exacte
\[
 \rho_9(n)=1+((n-1)\bmod9),\qquad
 q_9(n)=\left\lfloor\frac{n-1}{9}\right\rfloor,
 \qquad n=\rho_9(n)+9q_9(n).
\]
Une position relevée est
$z=(x,y,q_x,q_y)\in D_9^2\times\mathbb Z^2$, où
$X=x+9q_x$ et $Y=y+9q_y$. L'incrément entier $u$ de la première
coordonnée agit par
\[
 T^x_u(x,y,q_x,q_y)=
 \bigl(\rho_9(x+u),y,q_x+c_9(x;u),q_y\bigr),\qquad
 c_9(x;u)=\left\lfloor\frac{x+u-1}{9}\right\rfloor.
\]
La deuxième coordonnée se transporte de la même manière. Les avances
cardinales correspondent à $u=\pm1$ ; l'incrément nul conserve l'état.
Le cocycle satisfait
\[
 c_9(x;u+v)=c_9(x;u)+c_9(\rho_9(x+u);v).
\]
En effet, les deux membres sont le quotient de l'unique décomposition de
$x+u+v$ avec résidu dans $D_9$. Par conséquent,
$T^x_vT^x_u=T^x_{u+v}$ et $T^x_{-u}=(T^x_u)^{-1}$ ; les transports
des deux coordonnées commutent. L'état conserve le nombre de franchissements
du bord, même lorsque la position visible revient.

\paragraph{Feuillets relevés et mise à jour des registres.}
La réalisation entière est fixée en prolongeant les opérations arithmétiques
APP aux coordonnées transportées :
\begin{align*}
 \widetilde S(z)&=X+Y=x+y+9(q_x+q_y),\\
 \widetilde P(z)&=XY
 =xy+9(q_xy+q_yx)+81q_xq_y.
\end{align*}
À chaque position, le registre des deux feuillets est conservé :
\[
 \bigl(\rho_9(\widetilde S),q_9(\widetilde S),
       \rho_9(\widetilde P),q_9(\widetilde P)\bigr).
\]
Si $e:z\to z'$ est un lien et $T$ le feuillet sélectionné, son incrément
intégral s'obtient à partir du registre avant et après le transport :
\begin{equation}
 \delta_e\widetilde T=
 \rho_9(\widetilde T(z'))-\rho_9(\widetilde T(z))
 +9\bigl(q_9(\widetilde T(z'))-q_9(\widetilde T(z))\bigr).
 \label{eq:cruce-incremento-registro}
\end{equation}
Cette identité explicite la conservation : le saut du résidu et
l’accroissement du quotient reconstruisent ensemble la variation du feuillet.
La sélection de feuillet, le transport cardinal et la mise à jour du
registre sont des opérations distinctes. Dans cette coordinatisation, on compose d’abord la
sélection ordonnée de feuillets $(T_x,T_y)$ avec $T_d$ et avec la mise à jour
de leurs quotients ; la formule ne postule pas que tout le TPK soit réduit
à cette représentation locale.

\begin{proposition}[Courbure entière et transport de frontière]
\label{prop:cruce-stokes-integral}
Définissons
$\widetilde A_x=\delta_x\widetilde T_x$ et
$\widetilde A_y=\delta_y\widetilde T_y$, au moyen de
\eqref{eq:cruce-incremento-registro}, et soit
\[
 \widetilde F_{xy}(X,Y)=
 \widetilde A_x(X,Y)+\widetilde A_y(X+1,Y)
 -\widetilde A_x(X,Y+1)-\widetilde A_y(X,Y).
\]
Les choix homogènes ont une courbure nulle. Les choix mixtes
$(\widetilde S,\widetilde P)$ et
$(\widetilde P,\widetilde S)$ ont respectivement pour courbures les entiers
$+1$ et $-1$. Pour le bord positif d'un rectangle
$R_{a,b}$ de côtés entiers $a,b>0$,
\[
 \widetilde W_{SP}(\partial R_{a,b})=ab,\qquad
 \widetilde W_{PS}(\partial R_{a,b})=-ab.
\]
Le résultat est indépendant du point de base et du nombre de franchissements
des raccords nonadiques.
\end{proposition}

\begin{proof}
La reconstruction exacte des deux feuillets donne
\[
 \delta_x\widetilde S=\delta_y\widetilde S=1,
 \qquad\delta_x\widetilde P=Y,
 \qquad\delta_y\widetilde P=X.
\]
Ainsi, la connexion $SP$ est $(1,X)$, de courbure
$1+(X+1)-1-X=1$ ; la connexion $PS$ est $(Y,1)$, de courbure
$Y+1-(Y+1)-1=-1$. Pour un feuillet unique, les différences
mixtes s'annulent. La somme des courbures des $ab$ plaquettes annule chaque
lien intérieur avec son inverse et conserve les liens de frontière.
On obtient ainsi $\pm ab$ dans $\mathbb Z$. Les quotients transportés
garantissent que la même identité vaut lors du franchissement de tout raccord.
\end{proof}

L'accumulateur de frontière peut à son tour être conservé sous forme nonadique :
$W=r_W+9q_W$. Pour une contribution entière $w_e$, il est mis à jour par
\[
 r_W'=\rho_9(r_W+w_e),\qquad
 q_W'=q_W+\left\lfloor\frac{r_W+w_e-1}{9}\right\rfloor.
\]
La composition par liens fournit la même somme que son évaluation
entière directe. Dans cette coordinatisation, le zéro entier a pour coordonnées
$(r_W,q_W)=(9,-1)$ ; c’est une représentation arithmétique de l’accumulateur,
sans identification supplémentaire avec les différents zéros orientés du TRIT.

\paragraph{Trois opérations sur l'orientation.}
Soit $\gamma$ une route cardinale fermée. Son inversion temporelle
$\gamma^{-1}$ inverse l'ordre des liens et l'orientation de chacun ;
par additivité,
\[
 \widetilde W_{SP}(\gamma^{-1})=-\widetilde W_{SP}(\gamma).
\]
L'échange des feuillets satisfait également
\[
 \widetilde W_{PS}(\gamma)=-\widetilde W_{SP}(\gamma).
\]
Pour le vérifier sur toute route, observons que la somme des deux
connexions est l'incrément exact de
$\widetilde S+\widetilde P=X+Y+XY$ ; son intégrale sur une boucle est nulle.

Le demi-tour simultané des deux directions agit, d’autre part,
par $J_c(X,Y)=(2c_x-X,2c_y-Y)$, avec $2c\in\mathbb Z^2$,
de sorte qu’il préserve la coordinatisation du réseau. Pour un lacet, on a
\[
 \widetilde W_{SP}(J_c\gamma)=\widetilde W_{SP}(\gamma).
\]
En effet, la contribution horizontale fermée est nulle et la verticale
se transforme comme
$\sum (2c_x-X)(-\delta Y)=\sum X\delta Y$, car
$\sum\delta Y=0$. Un demi-tour spatial préserve donc
l’orientation aréale de cette coordinatisation. Sa combinaison avec l’inversion temporelle
ou avec l’échange de feuillets se calcule en composant les opérations
correspondantes.

La différence est pertinente pour le registre électronique : l'inversion
simultanée des deux déplacements tous les $27$ pas d'avancement produit, dans
des fenêtres de $9$, la suite
\[
 \tau_j=(-1)^{\lfloor j/3\rfloor},\qquad 0\leq j<12,
 \qquad (+,+,+,-,-,-,+,+,+,-,-,-).
\]
Cette suite enregistre les orientations d'avancement. L'holonomie requiert
en outre les liens parcourus, les feuillets sélectionnés et les quotients
accumulés. Ainsi, le transport de l'orientation électronique est conservé
sans identifier un demi-tour des deux axes à l'inversion temporelle
d'une boucle.

\paragraph{Application au croisement marqué et portée sectorielle.}
Pour le croisement rectangulaire de côtés $4$ et $7$, cette réalisation évalue
exactement
\[
 \widetilde W_{SP}=28=1+9\cdot3,\qquad
 \widetilde W_{SP}(\gamma^{-1})=-28=8+9(-4).
\]
L’entier conserve l’aire orientée ; le résidu isolé représente seulement
sa classe nonadique. Par exemple, à partir de $(7,1)$, les représentants
périodiques des liens ont pour somme $-35$, tandis que la correction de
retenue ajoute $63$ et retrouve $28$.

Dans l'incidence marquée du lecteur angulaire, le $4$ est
$|H_e\cap P_\pi|$ et le $7$ est le pivot marqué $p_\varphi$.
La lecture rectangulaire fournit une réalisation entière du produit
$|H_e\cap P_\pi|\,p_\varphi$ une fois ce croisement déclaré.
L'application qui identifie cette région au secteur $12$ doit conserver
la marque et l'orientation du système sectoriel de
la Définition~\ref{def:ii09-sistema-sectorial}. Les affectations
sectorielles $12,23,13$ et leurs combinaisons de $A,h,\Delta^\circ$
sont une spécification supplémentaire à la loi d'holonomie rectangulaire.
Le présent résultat matérialise le relèvement d'aire, la retenue arithmétique et
leurs involutions, avec ces données de départ explicites.


\begin{definition}[Système sectoriel d'incidence angulaire]
\label{def:ii09-sistema-sectorial}
Le système sectoriel se compose du sextuplet ordonné d'invariants
\[
 (b,p,q,h_6,o_8,t)=(4,5,7,6,8,3)
\]
et des applications linéaires, sur les coordonnées angulaires
\((a,u,d)\),
\begin{align}
 \mathscr A_{12}^{\rm ang}(a,u,d)&=2a-pu+bq\,d,\label{eq:ii09-sector12}\\
 \mathscr A_{23}^{\rm ang}(a,u,d)&=qu-\frac{p^2}{2}\,d,\label{eq:ii09-sector23}\\
 \mathscr A_{13}^{\rm ang}(a,u,d)&=a-pb\,u-\frac{o_8-h_6}{t}\,d.
 \label{eq:ii09-sector13}
\end{align}
Les symboles \(h_6,o_8\) conservent les cardinaux de l'incidence
marquée ; \(t=3\) est la normalisation tritique déclarée de ce secteur.
Le sextuplet et les trois règles forment conjointement le système de
départ de cette construction.
\end{definition}

\begin{theorem}[Évaluation exacte du système sectoriel]
\label{thm:ii09-matriz-sectorial}
L'application conjointe
\((\mathscr A_{12}^{\rm ang},\mathscr A_{23}^{\rm ang},
\mathscr A_{13}^{\rm ang})^{\mathsf T}\) de la
Définition~\ref{def:ii09-sistema-sectorial} a pour unique matrice
\[
 \mathsf M_{\rm CKM}=
 \begin{pmatrix}
 2&-5&28\\
 0&7&-25/2\\
 1&-20&-2/3
 \end{pmatrix}.
\]
Les trois colonnes sont les images des vecteurs de coordonnées
\((1,0,0)\), \((0,1,0)\) et \((0,0,1)\) par ces applications.
\end{theorem}
\begin{proof}
Les opérations sectorielles donnent
\[
 bq=4\cdot7=28,\qquad p^2/2=5^2/2=25/2,
 \qquad pb=5\cdot4=20,\qquad
 (o_8-h_6)/t=(8-6)/3=2/3.
\]
Les images successives de la base de coordonnées sont donc
\[
 c_A=(2,0,1)^{\mathsf T},\qquad
 c_h=(-5,7,-20)^{\mathsf T},\qquad
 c_\Delta=(28,-25/2,-2/3)^{\mathsf T}.
\]
Une application linéaire est déterminée par ses images sur une
base. La matrice dont les colonnes sont ces images est celle indiquée.
La démonstration utilise le système sectoriel complet de la définition.
\end{proof}

Sur les coordonnées précédemment produites, on pose
\[
 (a,u,d)=\left(A_{\deg},\frac{C^*_{\deg}}6,
                       \frac{180}{\pi}\Delta_4\right).
\]
Les contributions directe, conjuguée et torsionnelle à chaque angle
sont ainsi individualisées. En particulier, la différence
\(o_8-h_6=2\) est conservée dans le coefficient du défaut torsionnel
du secteur \(13\). Cette présence concrète permet de suivre
l'incidence jusqu'au lecteur angulaire, en maintenant distincts le
système sectoriel, son évaluation et la paramétrisation unitaire.
